\chapter{Advance Care Planning}
\index{Advance Care Planning}

\section*{Definition and Overview}
Advance care planning is the process of thinking about, talking about, and writing down your preferences for future health care, for a time when you might not be able to speak for yourself. It involves reflecting on your values and what matters most to you, choosing a trusted person to make medical decisions on your behalf if you cannot, and completing documents that record your wishes.

These documents are known by different names in different places, including advance care directive, advance health directive, living will, and enduring power of attorney for health and medical decisions. Advance care planning is not only for older people or people with terminal illness; it can be valuable at any age, particularly after a serious diagnosis, with a progressive or chronic illness, or before major surgery. The goal is not to plan for death so much as to make sure that future care -- whatever it involves -- reflects the person's own values.

\section*{Epidemiology}
Most people say that talking about end-of-life preferences is important, but far fewer have actually had the conversation or written anything down. Surveys in several countries consistently show that only a minority of adults have completed an advance care plan, even though the majority think it would be useful.

People who do complete a plan are more likely to receive care that matches their wishes, are less likely to undergo unwanted or burdensome treatment at the end of life, and their families report less distress, guilt, and conflict when decisions must be made. The terminology and legal requirements differ between countries and, in many countries, between states and territories, so documents must be prepared according to local law.

\section*{Aetiology and Pathophysiology}
Serious illness, a sudden accident, a stroke, or dementia can leave a person unable to communicate their wishes at the very time that difficult decisions must be made.

\begin{itemize}
    \item \textbf{Without a plan:} family members and clinicians must guess what the person would have wanted, which can cause distress, disagreement, and guilt
    \item \textbf{Default treatment:} life-sustaining treatment such as cardiopulmonary resuscitation, mechanical ventilation, or artificial feeding may be provided by default even when it does not match the person's values
    \item \textbf{With a plan:} a clear document lets clinicians deliver treatment that fits the person's expressed preferences rather than the standard approach
    \item \textbf{Decision-making:} a plan allows the person to decide in advance who they trust to make decisions for them, and sets out how they want that person to act
    \item \textbf{Practical wishes:} it can also record where the person would prefer to be cared for, such as at home rather than in hospital
\end{itemize}

\section*{Clinical Features}
A well-rounded advance care plan typically covers several elements, all of which reflect the person's values and preferences.

\begin{itemize}
    \item \textbf{Values and goals:} what matters most -- for example, comfort, independence, being with family, or spiritual beliefs
    \item \textbf{Treatment preferences:} wishes about resuscitation, ventilation, artificial feeding and fluids, and hospital admission
    \item \textbf{Substitute decision-maker:} the trusted person appointed to make health decisions if the person cannot
    \item \textbf{Place of care:} preferences about being cared for at home, in hospital, or in residential aged care
    \item \textbf{Other wishes:} organ and tissue donation preferences, which are handled separately in some jurisdictions, and personal requests such as who should be contacted
\end{itemize}

\section*{Differential Diagnosis}
The terms used in advance care planning are often confused, and they mean different things.

\begin{itemize}
    \item An \textbf{advance care directive} records a person's wishes and values about future treatment and is usually prepared with a doctor's involvement
    \item An \textbf{enduring power of attorney (medical)} appoints a person to make decisions, rather than recording treatment wishes
    \item A \textbf{do-not-resuscitate (DNR) order} is a clinical document, made with a doctor, that applies to specific resuscitation in specific circumstances
    \item A \textbf{last will and testament} deals with property and money after death and is a separate legal matter
\end{itemize}

\section*{When to Seek Medical Care}
Advance care planning is not something you see a doctor for in the way you would for an illness; it is a conversation to have while you are well enough to make decisions, not in the middle of a crisis.

\begin{itemize}
    \item Start at any age, and review the plan after major life events
    \item It is especially relevant after a serious diagnosis, with a progressive or chronic illness, or before major surgery
    \item Speak to a GP, who can help you think through the issues and can make sure your plan reaches hospital records
    \item Seek legal advice for documents that require formal witnessing or registration in your jurisdiction
\end{itemize}

\section*{Diagnosis and Investigations}
Completing advance care planning involves a series of practical steps rather than medical tests.

\begin{itemize}
    \item \textbf{Conversation:} discussing values, fears, and preferences with family and a trusted doctor or care professional
    \item \textbf{Choosing a decision-maker:} selecting a willing and trusted person who understands the role and is prepared to act on your behalf
    \item \textbf{Documentation:} completing the correct forms for your state, territory, or country, with the required signing and witnessing
    \item \textbf{Sharing:} giving copies to family, the GP, and other services, and registering the plan where a register exists
    \item \textbf{Review:} revisiting the plan periodically and after new diagnoses, changes in health, or the death of a decision-maker
\end{itemize}

\section*{Management}
A good plan is one that is written, shared, and reviewed, not simply discussed.

\begin{itemize}
    \item Choose a substitute decision-maker who is willing and able to speak for you
    \item Talk openly with family and your GP about your wishes, not just about the document
    \item Complete the legally valid forms for your jurisdiction
    \item Provide copies to your GP, family, and the hospital each time you are admitted
    \item Register the plan if health systems offer a register
    \item Review it regularly, especially after major life events
    \item Mention your plan whenever you are admitted to hospital or aged care
\end{itemize}

\section*{Complications}
Most problems arise from not having a plan, or from having one that is unclear or outdated.

\begin{itemize}
    \item Receiving unwanted, burdensome treatment because no one knows the person's wishes
    \item Receiving less treatment than the person would have wanted, when decisions are made by default
    \item Family conflict and guilt over difficult decisions made under pressure
    \item Distress for loved ones who must guess what the person would have wanted
    \item Invalid or outdated documents that do not reflect current wishes
    \item Care that is inconsistent with the person's values and quality of life
\end{itemize}

\section*{Prognosis and Outcomes}
Most people who complete advance care planning report reassurance and peace of mind, and their families report less anxiety, guilt, and conflict. People with a plan are more likely to receive care that matches their stated preferences and are less likely to receive treatments they would not have chosen. A plan does not guarantee a specific outcome, and clinical decisions are always guided by doctors at the time in light of the situation, but it ensures that the person's voice is heard when they cannot speak for themselves.

\section*{Prevention and Self-Care}
\begin{itemize}
    \item \textbf{Plan early:} start while you are well, rather than in a crisis
    \item \textbf{Choose carefully:} select your substitute decision-maker deliberately, and confirm they are willing
    \item \textbf{Keep talking:} have the conversation more than once, and after major life changes
    \item \textbf{Make it accessible:} keep your documents where people can find them, and tell them where they are
    \item \textbf{Update regularly:} revise the plan after a new diagnosis, a change in health, or a change in family circumstances
    \item \textbf{Bring it with you:} take a copy of your plan whenever you are admitted to hospital
    \item \textbf{Seek help:} ask a GP, aged-care professional, or community legal service if you are unsure how to begin
\end{itemize}

\section*{Sources and Further Reading}
The following sources provide reliable, up-to-date information on advance care planning.

\begin{itemize}
    \item NHS. \textit{End of life care: planning ahead and advance decisions.} https://www.nhs.uk/conditions/
    \item Mayo Clinic. \textit{Advance care planning: advance directives and end-of-life decisions.} https://www.mayoclinic.org
    \item MedlinePlus. \textit{Advance care directives: what they are and how to complete them.} https://medlineplus.gov
    \item National Institute on Aging. \textit{Advance care planning: getting started.} https://www.nia.nih.gov
    \item World Health Organization. \textit{Palliative care and end-of-life care guidance.} https://www.who.int
\end{itemize}
